\documentclass[10pt,a4paper]{amsart}
\usepackage[margin=19mm]{geometry}
\usepackage{amsmath,amssymb,mathtools}
\usepackage[scr=rsfs,cal=euler]{mathalfa}
\usepackage{xfrac}
\usepackage[unicode,hidelinks,bookmarks=false]{hyperref}
\newcommand{\ie}{i.e.\ }
\newcommand{\cf}{cf.\ }
\newcommand{\resp}{respectively\ }
\newcommand{\Subsection}{\subsection}
\providecommand{\nobreakdash}{\nobreak\hbox{-}}
\setlength{\emergencystretch}{2em}
\multlinegap0pt
\usepackage{amssymb}
\usepackage{mathtools}
\newtheorem{lemma}{Lemma}
\newcommand{\Sn}{\mathfrak{S}_n}
\newcommand{\norm}[1]{\left\lVert #1 \right\rVert}
\newcommand{\psdge}{\succeq}
\newcommand{\psdle}{\preceq}
\title{A Resolution of the SS--RS--GD Inequalities}
\author{Binghui Peng}
\date{Source: arXiv:2607.22620v1 (CC BY 4.0)}
\begin{document}
\maketitle

\noindent\emph{Source and licence.} A Resolution of the SS--RS--GD Inequalities. Version arXiv:2607.22620v1 (CC BY 4.0), distributed by arXiv under the Creative Commons Attribution 4.0 International licence, http://creativecommons.org/licenses/by/4.0/, as recorded in the arXiv licence metadata for this exact version and shown on the abstract page https://arxiv.org/abs/2607.22620v1. Full licence notice: \texttt{licenses/arxiv-2607.22620-CC-BY-4.0.txt}.\par\smallskip

\noindent\textit{Editorial scope.} Counted unit: the article's lemma lem:amgm (source0.tex lines 339--345). The article's complete original proof of it is delivered with it (source0.tex lines 347--416, one \texttt{proof} environment of 70 lines). No mathematics is written, repaired or re-derived: the statement and the proof are the article's own lines, in the article's own notation, and no retained line is substituted or re-flowed.  No further source-local context is needed: every cross-reference of every retained line resolves inside the two delivered spans.  Nothing was omitted from any retained span, so the package carries no omission record.  No retained line contains a citation, so the wrapper carries no bibliography and no bibliography entry.  No retained line contains a figure environment, an image-inclusion command or drawing code, so no image is delivered and the package declares no asset dependency.  Editorial formatting only, in addition: an AMS \texttt{amsart} preamble that restates the article's own declarations and macros used by the retained lines, namely the environments \texttt{lemma} on separate counters, together with the standard packages those lines need.  The retained environments print in this document as Lemma 1 in order of appearance, whereas the article prints the same items with the numbering of its own full document.  The complete native source of the article as distributed in the arXiv e-print for this exact version is bundled unchanged as \texttt{sources/205960/source0.tex}.
\begin{lemma}[Near-identity shuffled product]\label{lem:amgm}
Let $X_1,\dots,X_n$ be symmetric $d\times d$ matrices with $C_i:=I+X_i\psdge0$ and
$\frac1n\sum_iC_i\psdle I$, and suppose $\delta:=\max_i\norm{X_i}\le\frac1{4n^2}$. Then
\[
0\;\psdle\;\widetilde R:=\frac1{n!}\sum_{\sigma\in\Sn}\prod_{i=1}^n C_{\sigma(i)}\;\psdle\;I.
\]
\end{lemma}

\begin{proof}
\emph{Symmetry and expansion.} As $\Sn$ is closed under reversal and reversing a product transposes
it, $\widetilde R^\top=\widetilde R$. Expanding $\prod_i(I+X_{\sigma(i)})$ and averaging over $\sigma$,
\begin{equation}\label{eq:expand}
\widetilde R=I+\sum_{k=1}^n E_k,\qquad
E_k=\frac1{k!}\sum_{\substack{(i_1,\dots,i_k)\\ \text{distinct}}}X_{i_1}X_{i_2}\cdots X_{i_k}.
\end{equation}
Indeed, for fixed $k$ each size-$k$ subset of the $n$ factor slots contributes, after averaging,
the same term $\frac{(n-k)!}{n!}\sum_{\text{distinct}}X_{i_1}\cdots X_{i_k}$ (sum over distinct
ordered $k$-tuples); there are $\binom nk$ such subsets, and $\binom nk\frac{(n-k)!}{n!}=\frac1{k!}$,
giving~\eqref{eq:expand}.

\emph{Low-order terms.} Put
\[
T:=-\sum_iX_i,\qquad Q:=\sum_iX_i^2\psdge0.
\]
The hypothesis $\frac1n\sum_iC_i=I+\frac1n\sum_iX_i\psdle I$ is exactly $\sum_iX_i\psdle0$, i.e.\
$T\psdge0$. From~\eqref{eq:expand},
\[
E_1=\sum_iX_i=-T,\qquad
E_2=\frac12\sum_{i\ne j}X_iX_j=\frac12\Bigl[\bigl(\textstyle\sum_iX_i\bigr)^2-\sum_iX_i^2\Bigr]
=\frac12\bigl(T^2-Q\bigr).
\]
Hence
\begin{equation}\label{eq:IR}
I-\widetilde R=T+\frac12\bigl(Q-T^2\bigr)-\sum_{k=3}^n E_k.
\end{equation}

\emph{Upper bound $\widetilde R\psdle I$.} Fix a unit vector $v$ and write $t:=v^\top Tv\ge0$ and
$q:=v^\top Qv=\sum_i\norm{X_iv}^2\ge0$. Since $T\psdge0$ and $\norm{T}\le\sum_i\norm{X_i}\le n\delta$,
the operator inequality $T^2\psdle\norm{T}\,T$ gives $v^\top T^2v\le n\delta\,t$, so
\begin{equation}\label{eq:firstpart}
v^\top\Bigl(T-\frac12T^2\Bigr)v\;\ge\;\Bigl(1-\frac{n\delta}2\Bigr)t\;\ge\;0,
\end{equation}
using $n\delta\le\frac1{4n}\le\frac12$. For the higher-order terms, each word with $k\ge3$ obeys,
by Cauchy--Schwarz on the two end factors and $\norm{X_{i_j}}\le\delta$ on the middle ones,
\[
\bigl|v^\top X_{i_1}\cdots X_{i_k}v\bigr|
\le\norm{X_{i_1}v}\,\delta^{\,k-2}\,\norm{X_{i_k}v}
\le\frac{\delta^{\,k-2}}2\bigl(\norm{X_{i_1}v}^2+\norm{X_{i_k}v}^2\bigr).
\]
Summing over distinct ordered $k$-tuples: for each fixed first index there are
$(n-1)(n-2)\cdots(n-k+1)=:(n-1)_{k-1}$ choices of the rest, and likewise for the last index, so
$\sum_{\text{tuples}}\norm{X_{i_1}v}^2=(n-1)_{k-1}\,q$ and the same for the last factor. Therefore
\begin{equation}\label{eq:Ekbound}
\bigl|v^\top E_kv\bigr|\le\frac{(n-1)_{k-1}}{k!}\,\delta^{\,k-2}\,q,
\qquad
\sum_{k=3}^n\bigl|v^\top E_kv\bigr|\le q\sum_{k=3}^n\frac{(n-1)_{k-1}}{k!}\,\delta^{\,k-2}.
\end{equation}
Using $(n-1)_{k-1}\le n^{\,k-1}$ and $\delta\le\frac1{4n^2}$, the $k$-th summand is at most
$\frac1{k!}n^{k-1}\delta^{k-2}=\frac1{k!}\,n\,(n\delta)^{k-2}\le\frac1{k!}\,n\,(4n)^{-(k-2)}$; the
$k=3$ term is $\le\frac16\cdot n\cdot\frac1{4n}=\frac1{24}$ and the series is dominated by a
geometric tail, giving $\sum_{k\ge3}\frac{(n-1)_{k-1}}{k!}\delta^{k-2}<\frac14$ for all $n\ge2$
(the supremum over $n$ is $\frac1{24}+o(1)<\frac14$). Combining~\eqref{eq:IR},
\eqref{eq:firstpart}, and~\eqref{eq:Ekbound},
\[
v^\top(I-\widetilde R)v\;\ge\;\Bigl(1-\frac{n\delta}2\Bigr)t+\frac12q-\frac14q\;\ge\;\frac14q\;\ge\;0,
\]
so $\widetilde R\psdle I$.

\emph{Lower bound $\widetilde R\psdge0$.} From~\eqref{eq:expand}, each $E_k$ is an average of at most
$\frac{n!}{(n-k)!}$ words of norm $\le\delta^k$, scaled by $\frac1{k!}$, so
$\norm{E_k}\le\binom nk\delta^k$ and
\[
\norm{\widetilde R-I}\le\sum_{k=1}^n\binom nk\delta^k=(1+\delta)^n-1\le e^{n\delta}-1\le
e^{1/(4n)}-1<1.
\]
Thus every eigenvalue of $\widetilde R$ exceeds $1-(e^{1/(4n)}-1)=2-e^{1/(4n)}>0$, so
$\widetilde R\succ0$. Combined with the upper bound, $0\psdle\widetilde R\psdle I$.
\end{proof}

\end{document}
